\section{Conclusion}
\label{sec:conclusion}

We introduced the ``no costume theories'' rule as a quantitative criterion for evaluating the empirical content of consciousness theories.
Using three complementary metrics---\costumeratio, prediction agreement, and \costumesens---applied across five major theories via systematic proxy evaluation, we find:

\begin{itemize}[leftmargin=*,itemsep=2pt,topsep=2pt]
    \item \IIT is overwhelmingly a costume theory, with 92\% representational-only predictions and an 84-point substrate sensitivity ($d = 5.09$).
    \item \GNW, \HOT, \RPT, and \AST form a behaviorally cohesive cluster (65--87\% agreement) with costume ratios of 0.20--0.30.
    \item The primary source of inter-theory disagreement in consciousness science is \IIT; removing it from comparison yields substantial consensus.
    \item The perfect simulation scenario serves as a pure diagnostic: \IIT is the only theory that denies consciousness to a system behaviorally identical to a conscious human.
\end{itemize}

These findings do not settle whether \IIT is \emph{wrong}---a theory can be unfalsifiable and still true.
But they establish that \IIT is, by a wide quantitative margin, the least empirically grounded of the five theories we studied.
The no costume theories rule provides a principled basis for prioritizing theories that make behavioral predictions, both in advancing the science of consciousness and in evaluating the growing question of AI consciousness.

\para{Future work.}
Three directions are most pressing:
(1)~multi-model replication using Claude, Gemini, and open-weight models to test whether our findings are model-specific;
(2)~expert validation, where consciousness researchers evaluate a subset of the proxy classifications for accuracy; and
(3)~refined costume tests with explicitly stipulated behavioral and processing equivalence across substrate conditions.
More broadly, we hope the \costumeratio framework inspires similar analyses in other fields where theories proliferate without clear empirical adjudication.
